
\documentclass{article}
\usepackage{colortbl}
\usepackage{makecell}
\usepackage{multirow}
\usepackage{supertabular}

\begin{document}

\newcounter{utterance}

\centering \large Interaction Transcript for game `cladder', experiment `full\_v1.5\_default', episode 1649 with score{-}Qwen3.5{-}9B{-}sp.
\vspace{24pt}

{ \footnotesize  \setcounter{utterance}{1}
\setlength{\tabcolsep}{0pt}
\begin{supertabular}{c@{$\;$}|p{.15\linewidth}@{}p{.15\linewidth}p{.15\linewidth}p{.15\linewidth}p{.15\linewidth}p{.15\linewidth}}
    \# & \multicolumn{2}{c}{Player} && \multicolumn{2}{c}{Game Master} \\
    \hline

    \theutterance \stepcounter{utterance}  
    & & & \multicolumn{4}{p{0.6\linewidth}}{
        \cellcolor[rgb]{0.9,0.9,0.9}{
            \makecell[{{p{\linewidth}}}]{
                \texttt{\tiny{[P1$\langle$GM]}}
                \texttt{You are an expert in causal inference. The following question is not a typical commonsense query, but rather a meticulously designed question created by a professor specializing in causal inference, intended to assess the students' mastery of the course content.} \\
\texttt{Imagine a self{-}contained, hypothetical world with only the following conditions, and without any unmentioned factors or causal relationships: Season has a direct effect on sprinkler and weather. Weather has a direct effect on ground. Sprinkler has a direct effect on ground. The overall probability of rainy season is 87\%. The probability of dry season and wet ground is 7\%. The probability of rainy season and wet ground is 50\%. Is the chance of wet ground larger when observing rainy season?} \\
\texttt{Start your answer with 'Yes' or 'No', followed by additional reasoning or evidence to support your explanation.} \\
            }
        }
    }
     \\ \\

    \theutterance \stepcounter{utterance}  
    & \multicolumn{4}{p{0.6\linewidth}}{
        \cellcolor[rgb]{0.561,0.737,0.561}{
            \makecell[{{p{\linewidth}}}]{
                \texttt{\tiny{[P1$\rangle$GM]}}
                \texttt{Yes} \\
\\ 
\texttt{**Reasoning:**} \\
\\ 
\texttt{To determine if the chance of wet ground is larger when observing a rainy season, we must compare the conditional probability \$P(\textbackslash{}text\{Wet Ground\} \textbackslash{}mid \textbackslash{}text\{Rainy Season\})\$ with the marginal probability \$P(\textbackslash{}text\{Wet Ground\})\$.} \\
\\ 
\texttt{1.  **Calculate the Marginal Probability of Wet Ground (\$P(\textbackslash{}text\{Wet Ground\})\$):**} \\
\texttt{    The problem provides the joint probabilities involving the ground state and the season. We can sum the probabilities of all scenarios where the ground is wet:} \\
\texttt{    *   \$P(\textbackslash{}text\{Rainy Season\} \textbackslash{}cap \textbackslash{}text\{Wet Ground\}) = 50\textbackslash{}\%\$} \\
\texttt{    *   \$P(\textbackslash{}text\{Dry Season\} \textbackslash{}cap \textbackslash{}text\{Wet Ground\}) = 7\textbackslash{}\%\$} \\
\\ 
\texttt{    Therefore:} \\
\texttt{    \$\$P(\textbackslash{}text\{Wet Ground\}) = 50\textbackslash{}\% + 7\textbackslash{}\% = 57\textbackslash{}\%\$\$} \\
\\ 
\texttt{2.  **Calculate the Conditional Probability of Wet Ground given Rainy Season (\$P(\textbackslash{}text\{Wet Ground\} \textbackslash{}mid \textbackslash{}text\{Rainy Season\})\$):**} \\
\texttt{    Using the definition of conditional probability \$P(A \textbackslash{}mid B) = \textbackslash{}frac\{P(A \textbackslash{}cap B)\}\{P(B)\}\$:} \\
\texttt{    *   We know \$P(\textbackslash{}text\{Rainy Season\} \textbackslash{}cap \textbackslash{}text\{Wet Ground\}) = 50} \\
            }
        }
    }
    & & \\ \\

    \theutterance \stepcounter{utterance}  
    & & & \multicolumn{2}{p{0.3\linewidth}}{
        \cellcolor[rgb]{0.9,0.9,0.9}{
            \makecell[{{p{\linewidth}}}]{
                \texttt{\tiny{[GM$|$GM]}}
                \texttt{Yes} \\
            }
        }
    }
    & & \\ \\

    \theutterance \stepcounter{utterance}  
    & & & \multicolumn{2}{p{0.3\linewidth}}{
        \cellcolor[rgb]{0.9,0.9,0.9}{
            \makecell[{{p{\linewidth}}}]{
                \texttt{\tiny{[GM$|$GM]}}
                \texttt{yes} \\
            }
        }
    }
    & & \\ \\

    \theutterance \stepcounter{utterance}  
    & & & \multicolumn{2}{p{0.3\linewidth}}{
        \cellcolor[rgb]{0.9,0.9,0.9}{
            \makecell[{{p{\linewidth}}}]{
                \texttt{\tiny{[GM$|$GM]}}
                \texttt{game\_result = WIN} \\
            }
        }
    }
    & & \\ \\

\end{supertabular}
}

\end{document}
